\section{Hankel determinants and the moment problem}
\label{sec:hankel}

The source page states three determinant formulas:
\[
 \det[\Cat_{i+j}]_{0\le i,j<n}=1,
 \qquad
 \det[\Cat_{i+j+1}]_{0\le i,j<n}=1,                              \tag{6.1}\label{eq:hankel01}
\]
and
\[
 \det[\Cat_{i+j+2}]_{0\le i,j<n}=n+1.                           \tag{6.2}\label{eq:hankel2}
\]
They have a particularly economical constructive proof.

\subsection{The ballot triangles}

For $0\le k\le r$, define
\[
 E_{r,k}=\binom{2r}{r-k}-\binom{2r}{r-k-1},                      \tag{6.3}\label{eq:even-ballot}
\]
and
\[
 O_{r,k}=\binom{2r+1}{r-k}-\binom{2r+1}{r-k-1}.                  \tag{6.4}\label{eq:odd-ballot}
\]
Outside the displayed range put the value zero.  Reflection proves that
$E_{r,k}$ counts nonnegative paths of $2r$ unit up/down steps ending at height
$2k$, and $O_{r,k}$ counts such paths of length $2r+1$ ending at height
$2k+1$.  Both arrays are lower triangular in $(r,k)$ and have diagonal entry
$1$.

Cut a Dyck path of semilength $r+s$ after $2r$ steps.  Its height there is
$2k$ for a unique $k$; the prefix is counted by $E_{r,k}$ and the reversed
suffix by $E_{s,k}$.  Therefore
\[
 \Cat_{r+s}=\sum_{k\ge0}E_{r,k}E_{s,k}.                          \tag{6.5}\label{eq:even-gram}
\]
Cutting a Dyck path of semilength $r+s+1$ after $2r+1$ steps similarly gives
\[
 \Cat_{r+s+1}=\sum_{k\ge0}O_{r,k}O_{s,k}.                        \tag{6.6}\label{eq:odd-gram}
\]
The cut, reversal, and gluing operations are primitive-recursive inverse maps,
so these are finite combinatorial identities, not analytic moment arguments.

\subsection{The two unit determinants}

Let $E^{(n)}=(E_{r,k})_{0\le r,k<n}$ and similarly $O^{(n)}$.  Equations
\eqref{eq:even-gram}--\eqref{eq:odd-gram} say
\[
 [\Cat_{i+j}]=E^{(n)}(E^{(n)})^{\mathsf T},
 \qquad
 [\Cat_{i+j+1}]=O^{(n)}(O^{(n)})^{\mathsf T}.                    \tag{6.7}
\]
Both factors are unit lower triangular.  Determinant multiplicativity and
transpose invariance now give Equation~\eqref{eq:hankel01}.  In a minimal
formalization, multiplicativity is needed only for triangular matrices and
their transposes; it can also be proved by bounded Gaussian elimination using
integer row operations.

\subsection{The determinant starting at \(\Cat_2\)}

Let $M^{(n)}$ be the $n$ by $(n+1)$ matrix
\[
 M^{(n)}_{r,k}=E_{r+1,k},
 \qquad 0\le r<n,
 \quad 0\le k\le n.
\]
Cutting after $2r+2$ steps gives
\[
 [\Cat_{r+s+2}]_{0\le r,s<n}=M^{(n)}(M^{(n)})^{\mathsf T}.       \tag{6.8}\label{eq:Mgram}
\]
Cauchy--Binet expresses its determinant as the sum of the squares of the
$n+1$ maximal minors of $M^{(n)}$.  Every one of those minors is $1$.

\begin{lemma}[Maximal minors of the shifted ballot matrix]
\label{lem:maximal-minors}
Deleting any one column of $M^{(n)}$ leaves an $n$ by $n$ matrix of determinant
$1$.
\end{lemma}

\begin{proof}
Deleting column $0$ leaves columns $1,\ldots,n$.  The result is unit lower
triangular because $E_{r+1,k}=0$ for $k>r+1$ and
$E_{r+1,r+1}=1$, so this minor is $1$.

For every $r\ge1$,
\[
 \sum_{k=0}^{r}(-1)^k E_{r,k}=0.                                 \tag{6.9}\label{eq:alternating-ballot}
\]
To verify this, substitute Equation~\eqref{eq:even-ballot}, shift the second
finite sum, and use the symmetry of the binomial coefficients in the expansion
of $(1-1)^{2r}$.  Hence the alternating vector
$(1,-1,1,\ldots,(-1)^n)^{\mathsf T}$ spans the kernel of $M^{(n)}$; the rank is
$n$ because the minor just exhibited is nonzero.

The signed cofactor vector
\[
 \bigl((-1)^j\det M^{(n)}_{\widehat j}\bigr)_{j=0}^{n}
\]
is also in that one-dimensional kernel.  Comparing its first coordinate with
the known unit minor shows it equals the alternating vector.  Thus every
unsigned maximal minor is $1$.  This argument uses only the Laplace expansion
and cancellation in $\Z$, all available as finite primitive-recursive
operations.
\end{proof}

Applying Cauchy--Binet to Equation~\eqref{eq:Mgram} now yields
\[
 \det[\Cat_{i+j+2}]_{0\le i,j<n}
 =\sum_{j=0}^{n}1^2=n+1,
\]
which is Equation~\eqref{eq:hankel2}.

\subsection{The two determinant conditions characterize the sequence}

Suppose a sequence $a_0,a_1,\ldots$ satisfies, for every positive $n$,
\[
 \det[a_{i+j}]_{0\le i,j<n}=1,
 \qquad
 \det[a_{i+j+1}]_{0\le i,j<n}=1.                                \tag{6.10}\label{eq:characterizing-Hankel}
\]
The $1$ by $1$ cases give $a_0=a_1=1$.  Assume $a_0,\ldots,a_{2r-1}$ have
already been determined.  In the first determinant of size $r+1$, the newest
term $a_{2r}$ occurs only in the southeast corner.  Its cofactor is the
preceding unshifted determinant, equal to $1$.  Therefore the equation
``determinant $=1$'' is a linear equation with coefficient $1$ and uniquely
determines $a_{2r}$.  The shifted determinant of size $r+1$ similarly
determines $a_{2r+1}$, with cofactor $1$.

Quantifier-free induction therefore proves uniqueness.  Since the Catalan
sequence satisfies Equation~\eqref{eq:characterizing-Hankel}, the two
conditions together characterize it, exactly as stated on Wikipedia.

\subsection{Positive Hankel forms and compact moment representations}

The Gram factorizations give a finite integral-free positivity theorem.  For
any integer vector $z=(z_0,\ldots,z_{n-1})$,
\[
 \sum_{i,j<n}z_iz_j\Cat_{i+j}
 =\sum_{k<n}\left(\sum_{i<n}z_iE_{i,k}\right)^2\ge0.             \tag{6.11}
\]
The determinant $1$ shows positive definiteness over $\Q$.  The density in Equation~\eqref{eq:moment-density} gives the compatible
compactly supported Stieltjes representation on $[0,4]$:
\[
 \sum_{i,j}z_iz_j\Cat_{i+j}
 =\int_0^4\left(\sum_i z_ix^i\right)^2d\mu(x).
\]
After $x=4y$, the normalized matrix
$[\Cat_{i+j}/4^{i+j}]$ has a Hausdorff representation on $[0,1]$.
Thus the finite combinatorial factorization and the constructive moment code
prove the same positivity by independent routes.
